\documentclass[11pt]{article}
\usepackage[T1]{fontenc}
\usepackage{lmodern}
\usepackage{amsmath,amssymb,amsthm,mathtools}
\usepackage{microtype}
\usepackage[margin=0.92in]{geometry}
\usepackage{booktabs,longtable,array}
\usepackage{enumitem}
\usepackage{xcolor}
\usepackage{listings}
\usepackage{hyperref}
\usepackage{fancyhdr}
\usepackage{tcolorbox}
\usepackage{graphicx}

\definecolor{ink}{HTML}{172033}
\definecolor{accent}{HTML}{294C7A}
\definecolor{soft}{HTML}{EEF3F8}
\definecolor{warn}{HTML}{7A3D29}
\definecolor{codebg}{HTML}{F7F8FA}

\hypersetup{
  colorlinks=true,
  linkcolor=accent,
  urlcolor=accent,
  citecolor=accent,
  pdftitle={Crowned Kind-Polymorphic Forest Reflection},
  pdfauthor={OpenAI research package for BusyLean}
}

\pagestyle{fancy}
\fancyhf{}
\fancyhead[L]{\small Crowned kind-polymorphic forest reflection}
\fancyhead[R]{\small BusyLean Contender 11 proposal}
\fancyfoot[C]{\thepage}
\setlength{\headheight}{14pt}
\setlength{\parindent}{0pt}
\setlength{\parskip}{0.55em}
\setlist{leftmargin=*,itemsep=0.25em,topsep=0.35em}
\emergencystretch=2em

\lstdefinelanguage{lean}{
  keywords={def,theorem,inductive,structure,where,match,with,fun,let,in,
    namespace,end,open,abbrev,import,by,have,show,exact,termination_by,
    cases,induction,obtain,refine,simpa,unfold,change,rw},
  sensitive=true,
  morecomment=[l]{--},
  morecomment=[s]{/-}{-/},
  morestring=[b]",
}
\lstset{
  language=lean,
  basicstyle=\ttfamily\footnotesize,
  keywordstyle=\bfseries\color{accent},
  commentstyle=\itshape\color{black!58},
  stringstyle=\color{warn},
  backgroundcolor=\color{codebg},
  columns=fullflexible,
  keepspaces=true,
  frame=single,
  framerule=0.3pt,
  rulecolor=\color{black!25},
  xleftmargin=0.5em,
  xrightmargin=0.5em,
  breaklines=true,
  breakatwhitespace=true,
  showstringspaces=false,
  tabsize=2,
  literate={→}{{$\to$}}1 {←}{{$\leftarrow$}}1 {∀}{{$\forall$}}1
    {∃}{{$\exists$}}1 {≤}{{$\le$}}1 {≥}{{$\ge$}}1 {≠}{{$\ne$}}1
    {∈}{{$\in$}}1 {×}{{$\times$}}1 {ω}{{$\omega$}}1
    {⟨}{{$\langle$}}1 {⟩}{{$\rangle$}}1 {·}{{$\cdot$}}1
}

\newtheorem{theorem}{Theorem}[section]
\newtheorem{lemma}[theorem]{Lemma}
\newtheorem{corollary}[theorem]{Corollary}
\theoremstyle{definition}
\newtheorem{definition}[theorem]{Definition}
\newtheorem{proposition}[theorem]{Proposition}
\theoremstyle{remark}
\newtheorem{remark}[theorem]{Remark}

\newcommand{\N}{\mathbb{N}}
\newcommand{\Fk}{\mathrm{F}_{\omega}^{\mathsf{k}}}
\newcommand{\code}[1]{\texttt{#1}}
\newcommand{\R}{\mathcal{R}_{11}}
\newcommand{\Sstage}{\mathcal{S}^{(11)}}
\newcommand{\Crown}{\mathsf{Crown}}
\newcommand{\FamTen}{\mathsf{family}_{10}}
\newcommand{\Feleven}{F_{11}}
\newcommand{\ContEleven}{C_{11}}
\newcommand{\ContTen}{C_{10}}
\newcommand{\repo}{\textsc{BusyLean}}
\newcommand{\pin}{\texttt{4c4b03d0e2003c479e3bee90a21c8bea13f66c40}}

\newtcolorbox{statusbox}{
  colback=soft,
  colframe=accent,
  boxrule=0.7pt,
  arc=2pt,
  left=7pt,right=7pt,top=6pt,bottom=6pt
}
\newtcolorbox{warningbox}{
  colback=warn!5,
  colframe=warn,
  boxrule=0.7pt,
  arc=2pt,
  left=7pt,right=7pt,top=6pt,bottom=6pt
}

\title{\textbf{Crowned Kind-Polymorphic Forest Reflection}\\[0.35em]
  \large A proposed Contender 11 for BusyLean}
\author{OpenAI research package prepared for the BusyLean repository}
\date{October 2026}

\begin{document}
\maketitle

\begin{abstract}
At the pinned revision \pin, the current \repo{} entry is Contender~10,
obtained by reflecting a kind-polymorphic extension of System~$F_\omega$
over Contender~9's numerical family.  This article proposes a structurally
stronger, still constructive entry.  First, the unary oracle of Contender~10
is lifted back to a countable family by the repository's certified $2$-adic
pairing.  Second, the resulting kind-polymorphic reflector is installed at
every node of a well-founded $W$-tree.  Third, the forest is run along the
repository's self-generated tree index, which joins every $W\N$ tree denoted
by the earlier W-language.  Finally, an input-diagonal ``crown'' of $n+1$
additional reflection nodes is placed above that index.

The proposed family is
\[
  \FamTen(n)=F_{10}(F_{10}(n)),\qquad
  \Feleven(n)=
  \Sstage_{\FamTen}\bigl(\Crown_{n+1}(I),n\bigr),
\]
and the proposed closed number is
\[
  \ContEleven=\Feleven(\Feleven(42)).
\]
The comparison is exact and needs no ordinal calibration:
$\FamTen(n)<\Feleven(n)$ for every $n$, hence
$\ContTen=\FamTen(42)<\ContEleven$.  Stronger structural statements show that
every crown edge is strict, every accepted depth-bounded kind-polymorphic
program at the root is exceeded, and every tree denoted by the old W-language
is eventually dominated.

The package supplies Lean modules, an axiom/dependency audit, integration
tooling, a finite independent reference model, and reproducibility evidence.
The packaging environment did not contain Lean~4.34.0 or \code{lake}; therefore
the formal files are a concrete kernel-checkable proposal, not a claim of a
completed kernel build.  Static auditing and 670 finite structural checks pass;
the remaining certification and timing obligations are stated explicitly.
\end{abstract}

\begin{statusbox}
\textbf{Certification status.}
The mathematical comparison and the proposed Lean proof are given in full.
The PDF, source audit, package hashes, and finite reference tests were produced
in the present environment.  Lean elaboration, \code{\#print axioms}, transitive
closure checking, and the 15/60 second budgets remain to be run in a pinned
Lean~4.34.0 checkout.  Nothing in this article should be read as silently
substituting Python tests for the Lean kernel.
\end{statusbox}

\tableofcontents

%=====================================================================
\section{Research target}

\subsection{The actual baseline}

The target is not the original Coq Contender~5 and not one of the earlier
unmerged proposals.  The inspected BusyLean revision is \pin.  At that
revision, the repository defines
\begin{lstlisting}
def family9 (n : Nat) : Nat := F9 (F9 n)
def F10 (n : Nat) : Nat := Fk.reflectFk family9 n
def contender_10 : Nat := F10 (F10 42)
\end{lstlisting}
and proves \code{contender\_9\_lt\_contender\_10}.  The reflected object
language is a kind-polymorphic extension of System~$F_\omega$ with a total
normalizer justified by reducibility candidates.  A bounded syntax list and a
sound checker make its finite numerical maximum a computable natural number.

Contender~10 is already a substantial entry.  A serious successor must not
merely add one, multiply by a fixed constant, or repeat the named scalar.  The
new construction below changes the 
\emph{shape of reflection}: the stronger object language is reflected at every
node of a well-founded forest, and the forest is traversed through a
self-generated index plus a diagonal crown.

\subsection{Contest obligations}

The repository rules impose five constraints that matter directly here.

\begin{enumerate}
\item The result must be a closed computable \code{Nat}.  Uncomputable Busy
  Beaver functions, \code{noncomputable}, \code{partial}, \code{unsafe},
  external implementations, and opaque placeholders are excluded.
\item The comparison with the preceding entry must be a formal strict
  inequality.
\item The permitted axiom envelope is the one already used by BusyLean:
  \code{propext} and \code{Quot.sound}; classical choice and custom axioms are
  excluded.
\item The new definition and proof phases have separate 15-second and
  60-second budgets once the accepted predecessor is precompiled.
\item The named scalar \code{contender\_10} may not be used as a shortcut.
  Reusing the preceding \emph{construction} as an oracle family is permitted
  when it is embedded in a qualitatively stronger construction.
\end{enumerate}

The proposed source is organized around these constraints.  The definition
module imports only the previous definition module.  It defines the old
construction uniformly as \code{family10 n := F10 (F10 n)}, but its executable
syntax contains no occurrence of \code{contender\_10}.  The proof and audit are
separate modules.

\subsection{What ``significantly exceeds'' means here}

A single scalar inequality is necessary but not, by itself, persuasive.  We
therefore use four exact witnesses of structural separation.

\begin{center}
\begin{tabular}{p{0.27\linewidth}p{0.64\linewidth}}
\toprule
Witness & Exact content \\
\midrule
Uniform family bound & $\FamTen(n)<\Feleven(n)$ for every $n$. \\
Crown ladder & Every consecutive crown level is strictly larger at every
input.  At input 42 there are 43 such edges; at the outer input there are
$\Feleven(42)+1$ edges. \\
Program reflection & Every accepted $\Fk$ program of depth at most $n$,
relative to the preceding crown level, is below $\Feleven(n)$. \\
Tree cofinality & Every $W\N$ tree denoted by the earlier W-language is
strictly below the new family from an explicit threshold onward. \\
\bottomrule
\end{tabular}
\end{center}

None of these statements relies on a conjectural ordinal label.

%=====================================================================
\section{The inherited semantic interfaces}

\subsection{Kind-polymorphic reflection}

Write $V_f(n,t)$ for \code{Fk.fvalue f n t}: type-check $t$ with fuel $n$;
if its type is the Church numeral type, erase and normalize it relative to the
unary oracle $f$, then decode the normal form; otherwise return zero.  Let
$T_n$ be the finite list \code{Fk.ctmsUpTo n}.  Contender~10 uses
\[
  R_{10}(f,n)=n+1+
  \max\left(f(n),\max_{t\in T_n}V_f(n,t)\right).
\]
The repository proves three facts that are crucial for reuse:

\begin{enumerate}
\item the checker is sound;
\item checked terms erase to strongly normalizing terms;
\item $R_{10}(f,n)$ exceeds both $f(n)$ and every listed program value.
\end{enumerate}

The proposed successor does not alter the object language, checker, normalizer,
or decoder.  It changes the oracle interface from one function to a countable
family and then iterates that reflection through a $W$-tree.

\subsection{Packing a family into one unary oracle}

For $k>0$, write uniquely
\[
  k=2^i(2m+1).
\]
The repository's functions \code{v2}, \code{unpair2}, and \code{packFamily}
implement
\[
  \operatorname{code}(i,m)=2^i(2m+1),\qquad
  \operatorname{pack}(f)(k)=f(\operatorname{unpair}_2(k)).
\]
It proves the exact decoder theorem
\[
  \operatorname{unpair}_2(2^i(2m+1))=(i,m),
\]
so
\[
  \operatorname{pack}(f)(2^i(2m+1))=f(i,m).
\]
This matters because a kind-polymorphic program at a forest node can call any
child $i$ at any numerical input $m$ using an ordinary unary oracle call.  No
uncomputable family oracle is introduced.

\subsection{The self-generated W-tree index}

The earlier W-language gives a total interpretation of raw terms into values
that include $W\N$ trees.  BusyLean defines
\[
\begin{aligned}
  \operatorname{defTree}(t) &:= \text{the }W\N\text{ tree denoted by }t,\\
  \operatorname{treesOfDepth}(k)
    &:= \operatorname{node}\bigl(i\mapsto
       \operatorname{defTree}(\operatorname{nthD}(\operatorname{allTerms}(k),i))\bigr),\\
  I &:= \operatorname{node}(\operatorname{treesOfDepth}).
\end{aligned}
\]
Thus child $k$ of $I$ joins every tree in the depth-$k$ enumeration.  The first
enumeration is empty, so its default term denotes a leaf.  This gives an exact
three-step route
\[
  \text{seed value}
  < \Sstage(\mathsf{leaf})
  < \Sstage(\operatorname{treesOfDepth}(0))
  < \Sstage(I),
\]
which will be the core of the comparison.

%=====================================================================
\section{The new family-valued reflector}

\begin{definition}[Family reflection]
For a family $f:\N\to\N\to\N$, define
\[
  \R(f,n)=n+1+
  \max\left(
    \max_{i\le n} f(i,n),
    \max_{t\in T_n} V_{\operatorname{pack}(f)}(n,t)
  \right).
\]
In Lean this is \code{reflect11}.
\end{definition}

\begin{lstlisting}
def reflect11 (fs : W.Family) (n : Nat) : Nat :=
  n + 1 + BigNumber.maxN (W.prefixMax (fun i => fs i n) n)
    (BigNumber.maximum (Fk.fvalue (packFamily fs) n)
      (Fk.ctmsUpTo n))
\end{lstlisting}

Every operation in this expression is total.  The two maxima are over finite
lists or finite prefixes.  The only potentially subtle component is
\code{Fk.fvalue}, whose totality is inherited from the formal strong
normalization theorem.

\begin{lemma}[Visible child bound]
If $i\le n$, then $f(i,n)<\R(f,n)$.
\end{lemma}
\begin{proof}
The prefix maximum is at least $f(i,n)$.  It is at most the outer maximum, and
$\R$ adds $n+1$, which is positive.
\end{proof}

\begin{lemma}[Program bound]
If $t\in T_n$, then
\[
  V_{\operatorname{pack}(f)}(n,t)<\R(f,n).
\]
\end{lemma}
\begin{proof}
The listed program value is at most the finite program maximum, hence at most
the outer maximum.  Again the positive summand makes the inequality strict.
\end{proof}

\begin{lemma}[Input bound]
For every $f,n$, $n<\R(f,n)$.
\end{lemma}
\begin{proof}
$\R(f,n)=n+1+M$ for a natural number $M$.
\end{proof}

These lemmas are deliberately order-free in the oracle.  Reflection is not
monotone in an arbitrary oracle: a program may inspect an oracle result and
branch in an anti-monotone way.  All comparisons in this package therefore use
\emph{explicit inclusion}---a child occurs in the finite prefix, or a program
occurs in the finite syntax list---rather than a false monotonicity principle.

%=====================================================================
\section{Installing the reflector on a well-founded forest}

\begin{definition}[Forest stage]
For a seed family $s$, define by structural recursion on a $W\N$ tree
\[
\begin{aligned}
  \Sstage_s(\mathsf{leaf},n) &= \R(s,n),\\
  \Sstage_s(\mathsf{node}(g),n)
    &=\R\bigl(i\mapsto \Sstage_s(g(i),\cdot),n\bigr).
\end{aligned}
\]
\end{definition}

\begin{lstlisting}
def stage11 (seed : W.Family) : BigNumber.WTree Nat → Nat → Nat
  | .leaf => reflect11 seed
  | .node g => reflect11 (fun i => stage11 seed (g i))
\end{lstlisting}

This is the same well-founded recursion pattern already accepted for the
repository's Contenders~6 and~9.  At a node, the unary object-language oracle is
the packed family of all child stages.  A program can choose a child and an
input dynamically.

\begin{lemma}[Strict parent edge]
If $i\le n$, then
\[
  \Sstage_s(g(i),n)<\Sstage_s(\mathsf{node}(g),n).
\]
\end{lemma}
\begin{proof}
At the parent, the oracle family is $j\mapsto\Sstage_s(g(j),\cdot)$.  Apply the
visible child bound at index $i$.
\end{proof}

\begin{corollary}[Every stage exceeds its input]
For every seed $s$, tree $a$, and input $n$,
\[
  n<\Sstage_s(a,n).
\]
\end{corollary}
\begin{proof}
Case split on $a$ and apply the input bound for $\R$.
\end{proof}

The parent-edge theorem is the engine behind both the self-generated index and
the crown.  It is unusually robust: no semantic comparison between two
oracles, and no analysis of program behavior, is required.

%=====================================================================
\section{The diagonal crown}

\subsection{Definition}

For a tree $a$, put
\[
  \Crown_0(a)=a,\qquad
  \Crown_{k+1}(a)=\mathsf{node}(i\mapsto\Crown_k(a)).
\]
Each crown node repeats one child at every index.  In particular, index zero is
visible at every input.

\begin{lstlisting}
def crown11 : Nat → BigNumber.WTree Nat → BigNumber.WTree Nat
  | 0, a => a
  | k + 1, a => .node fun _ => crown11 k a
\end{lstlisting}

\begin{lemma}[Crown step]
For all $s,a,k,n$,
\[
  \Sstage_s(\Crown_k(a),n)
  <\Sstage_s(\Crown_{k+1}(a),n).
\]
\end{lemma}
\begin{proof}
The left-hand tree is child zero of the right-hand node.  Since $0\le n$, the
strict parent-edge theorem applies.
\end{proof}

\begin{corollary}[Positive crown]
For all $k,n$,
\[
  \Sstage_s(a,n)<\Sstage_s(\Crown_{k+1}(a),n).
\]
\end{corollary}
\begin{proof}
Chain $k+1$ crown-step inequalities.  The supplied Lean proof performs this
chaining by induction on $k$.
\end{proof}

\subsection{Why the height is input-diagonal}

A fixed crown of height $43$ would strengthen the scalar at input 42 but would
not define a uniformly expanding family.  We instead use height $n+1$.  This
has three benefits.

\begin{enumerate}
\item The uniform theorem records a growing number of explicit semantic
reflection levels, not a hard-coded scalar embellishment.
\item At the advertised inner input 42, there are exactly 43 strict crown
edges above the self-generated index.
\item At the outer input $\Feleven(42)$, the closed contender traverses
$\Feleven(42)+1$ fresh crown edges.  This is a diagonal amplification whose
height is itself already the new family value.
\end{enumerate}

The crown is not the only source of growth: beneath it lies the entire
self-generated forest.  Its purpose is to provide a clean, exact, inspectable
margin in addition to the more complex forest semantics.

%=====================================================================
\section{The proposed family and closed number}

Define the predecessor uniformly by
\[
  \FamTen(n)=F_{10}(F_{10}(n)).
\]
By definitional equality, $\FamTen(42)=\ContTen$.  Let the seed family ignore
its child index and return $\FamTen$:
\[
  s(i,n)=\FamTen(n).
\]
Then define
\[
  \Feleven(n)=\Sstage_s(\Crown_{n+1}(I),n),
  \qquad
  \ContEleven=\Feleven(\Feleven(42)).
\]

\begin{lstlisting}
def family10 (n : Nat) : Nat := F10 (F10 n)

def F11 (n : Nat) : Nat :=
  stage11 (fun _ => family10) (crown11 (n + 1) index) n

def contender_11 : Nat := F11 (F11 42)
\end{lstlisting}

The definition module never names \code{contender\_10}.  It reuses
\code{F10}, just as earlier accepted BusyLean entries reuse the previous
uniform construction as an oracle.  The new scalar is therefore not
``the old scalar plus a margin''; it is a new reflection hierarchy whose seed
happens to contain the predecessor family.

%=====================================================================
\section{Exact comparison with Contender 10}

\subsection{The index already beats the seed}

Let $s(i,n)=\FamTen(n)$.  At every input $n$:

\begin{enumerate}
\item The leaf stage is $\R(s,n)$, and child zero of the seed prefix is
$s(0,n)=\FamTen(n)$.  Hence
\[
  \FamTen(n)<\Sstage_s(\mathsf{leaf},n).
\]
\item The first child of \code{treesOfDepth 0} is the denotation of the default
zero term, which the repository proves is a leaf.  Therefore
\[
  \Sstage_s(\mathsf{leaf},n)
  <\Sstage_s(\operatorname{treesOfDepth}(0),n).
\]
\item \code{treesOfDepth 0} is child zero of $I$, so
\[
  \Sstage_s(\operatorname{treesOfDepth}(0),n)
  <\Sstage_s(I,n).
\]
\end{enumerate}

Combining these inequalities gives the formal theorem
\[
  \FamTen(n)<\Sstage_s(I,n).
\]
This proof is intentionally concrete.  It does not need to know the height or
ordinal represented by $I$.

\subsection{Adding the crown}

By the positive-crown theorem with $k=n$,
\[
  \Sstage_s(I,n)<\Sstage_s(\Crown_{n+1}(I),n)=\Feleven(n).
\]
Hence:

\begin{theorem}[Uniform predecessor bound]
For every $n\in\N$,
\[
  \FamTen(n)<\Feleven(n).
\]
The supplied Lean name is \code{family10\_lt\_F11}.
\end{theorem}

\subsection{Closing the diagonal}

At $n=42$, definitional equality gives
\[
  \ContTen=\FamTen(42)<\Feleven(42).
\]
Every forest stage exceeds its input, so at the outer stage
\[
  \Feleven(42)
  <\Sstage_s(\Crown_{\Feleven(42)+1}(I),\Feleven(42))
  =\Feleven(\Feleven(42))=\ContEleven.
\]
Therefore:

\begin{theorem}[Required scalar comparison]
\[
  \ContTen<\ContEleven.
\]
The supplied Lean name is
\code{contender\_10\_lt\_contender\_11}.
\end{theorem}

The proof does not unfold a single enormous numeral.  It composes the uniform
family inequality with the general input bound.

%=====================================================================
\section{Stronger separation theorems}

\subsection{Root-program reflection}

Because
\[
  \Crown_{n+1}(I)=
  \mathsf{node}(i\mapsto\Crown_n(I)),
\]
the root oracle of $\Feleven(n)$ is the constant family
\[
  i\mapsto \Sstage_s(\Crown_n(I),\cdot).
\]
Consequently every accepted program $t\in T_n$ satisfies
\[
V_{\operatorname{pack}
  (i\mapsto \Sstage_s(\Crown_n(I),\cdot))}(n,t)<\Feleven(n).
\]
This is exactly \code{root\_program\_lt\_F11}.  It says more than
``the previous value occurs in a maximum'': the entire finite fragment of the
strong object language, relative to the entire preceding crown level, is
reflected at the root.

\subsection{Eventual domination of every old-language tree}

Suppose a raw W-language term $t$ belongs to the depth-$k$ enumeration, and let
$i$ be an index where it occurs.  For every
$n\ge k+i$, both $k\le n$ and $i\le n$.  The denoted tree is then a visible
child of \code{treesOfDepth k}; that join is a visible child of $I$; and $I$ is
strictly below the crown.  Thus
\[
  \Sstage_s(\operatorname{defTree}(t),n)<\Feleven(n).
\]
The threshold $k+i$ is explicit.  The Lean theorem is
\code{defined\_tree\_eventually\_lt11}.

This cofinality result is relative to the earlier W-language's tree output.  It
does not claim that kind-polymorphic programs themselves generate every tree;
the current $\Fk$ object language returns Church numerals, not $W$-objects.
That distinction is a natural target for a future, stronger contender.

\subsection{A strict semantic ladder}

For fixed $n$, set
\[
  A_k(n)=\Sstage_s(\Crown_k(I),n).
\]
Then
\[
  A_0(n)<A_1(n)<\cdots<A_{n+1}(n)=\Feleven(n).
\]
At $n=42$, the package therefore exhibits a 43-edge strict chain above the
already self-generated index.  At the final outer input, it exhibits a chain
of length $\Feleven(42)+1$.  This is a structural count of reflection nodes,
not a claim that successive numerical differences are one; each edge encloses
a new finite maximum over object-language computations relative to the lower
level.

%=====================================================================
\section{Constructivity and totality}

The proposed number is computable in principle for the same reason as the
accepted predecessors, plus two elementary structural recursions.

\begin{enumerate}
\item \textbf{Finite enumeration.}  \code{Fk.ctmsUpTo n} is a finite list
constructed recursively from finite lists of kinds, types, and terms.
\item \textbf{Decidable checking.}  \code{Fk.infer} is fuel bounded.  Rejected
terms contribute zero.
\item \textbf{Total normalization.}  An accepted term at the Church numeral
type has a proof of strong normalization, so \code{Fk.nf} is a total Lean
function on the proof-carrying term.
\item \textbf{Finite maxima.}  Both the syntax maximum and the child-prefix
maximum range over finite data.
\item \textbf{Well-founded forest recursion.}  \code{stage11} recurses
structurally on \code{WTree}, exactly as the existing \code{stage9} does.
\item \textbf{Primitive crown recursion.}  \code{crown11} recurses on a natural
number and constructs a finite unary spine.
\item \textbf{Closed specialization.}  \code{contender\_11} substitutes the
closed numeral 42 into closed total functions.
\end{enumerate}

No semantic oracle is assumed.  The only oracle seen by an object-language
program is a Lean-computable function generated from recursive calls on
strict subtrees.  The result is fantastically infeasible, but infeasibility is
not noncomputability.

\begin{warningbox}
\textbf{Never evaluate the actual contender.}
Computability is certified intensionally by total definitions and termination
proofs.  Attempting to normalize \code{F11 42} would be both unnecessary and
catastrophic.  The comparison proofs quantify over symbolic inputs and use
finite-maximum inequalities.
\end{warningbox}

%=====================================================================
\section{Why this is a qualitative advance}

\subsection{From one reflection to a reflection forest}

Contender~10 applies the kind-polymorphic reflector once to \code{family9} at
each input.  The new construction makes that reflector the local transition
at \emph{every node} of a well-founded tree.  At a node, an object-language
program has uniform access to all child stages through the packed oracle.  The
node then exceeds all visible children and all accepted programs relative to
them.

This is not equivalent to a fixed iteration count written in arithmetic.  The
shape of the iteration is a $W$-tree whose branching can be selected by
computed indices.  The self-generated index ranges over every tree output by
the earlier W-language, and the crown adds a separate diagonal sequence of
new layers.

\subsection{What is proved, and what is heuristic}

The following claims are exact consequences of the supplied definitions and
proof scripts:

\begin{itemize}
\item $\FamTen(n)<\Feleven(n)$ for every $n$;
\item $\ContTen<\ContEleven$;
\item each crown edge is strict;
\item all root programs in the bounded $\Fk$ enumeration are reflected;
\item every W-language-defined tree is eventually dominated.
\end{itemize}

It is reasonable to describe the move from a single $\Fk$ reflection to
$\Fk$ reflection along a self-generated forest as a multiplicative or
iteration-level increase in proof-theoretic terms.  However, this article does
\emph{not} assign an exact ordinal to the kind-polymorphic language, to the
forest closure, or to the crowned diagonal.  The repository's formal
inequalities are stronger evidence than an unsupported ordinal name.

\subsection{No claim of maximality}

No computable description can be the last word in an unrestricted contest: a
new description can reflect or diagonalize over it.  The research objective is
therefore a transparent and formally comparable advance over the actual
incumbent, not a metaphysical ``largest computable number.''

%=====================================================================
\section{No-laziness and dependency hygiene}

The computational dependency graph is designed to satisfy the repository's
no-laziness rule.

\begin{enumerate}
\item \code{Contender11.lean} imports \code{BusyLean.Contender10}, not
\code{BusyLean.Comparison10}.
\item Its executable syntax contains no identifier \code{contender\_10}.
\item The predecessor enters only through
\code{family10 n := F10 (F10 n)}, a uniform family subsequently used as the
leaf seed of a new forest operator.
\item The new scalar applies \code{F11} to its own inner value, rather than
adding a large but semantically inert arithmetic constant.
\end{enumerate}

The package includes a kernel-side transitive closure command modeled on the
repository's audit.  It is intended to establish that
\code{BusyLean.contender\_11} avoids every earlier named scalar contender,
while still exposing all legitimate dependencies on the generic constructions.
A separate static audit strips nested Lean comments and strings before checking
that \code{contender\_10} is absent from executable definition syntax.

%=====================================================================
\section{Formalization inventory}

\begin{longtable}{p{0.34\linewidth}p{0.58\linewidth}}
\toprule
Lean declaration & Role \\
\midrule
\endfirsthead
\toprule
Lean declaration & Role \\
\midrule
\endhead
\code{reflect11} & Family-valued $\Fk$ reflection step. \\
\code{stage11} & Well-founded forest recursion installing the reflector at
all nodes. \\
\code{crown11} & Finite unary spine above a base tree. \\
\code{family10} & Uniform predecessor construction; definitionally equals
Contender 10 at input 42. \\
\code{F11} & Crowned forest family. \\
\code{contender\_11} & Closed proposed number. \\
\code{oracle\_lt\_reflect11} & Visible child bound. \\
\code{fvalue\_lt\_reflect11} & Bounded-program bound. \\
\code{input\_lt\_stage11} & Every stage exceeds its input. \\
\code{stage11\_crown\_step\_lt} & Every crown edge is strict. \\
\code{family10\_lt\_stage11\_index} & Explicit leaf--join--index comparison. \\
\code{family10\_lt\_F11} & Uniform predecessor-family comparison. \\
\code{defined\_tree\_eventually\_lt11} & Cofinality over old-language trees. \\
\code{root\_program\_lt\_F11} & Root reflects programs relative to the
preceding crown. \\
\code{contender\_10\_lt\_contender\_11} & Required scalar comparison. \\
\bottomrule
\end{longtable}

The proof module is intentionally short relative to the inherited
normalization development.  Its substantive work is orchestration: existing
soundness and normalization justify \code{Fk.fvalue}; existing pairing exposes
child families; new finite-maximum and tree-edge arguments establish the
hierarchy.

%=====================================================================
\section{Validation performed for this package}

\subsection{Static Lean-source audit}

The static audit passed all advertised checks.  It verifies, after removing
nested comments and strings:

\begin{itemize}
\item absence of \code{noncomputable}, \code{partial}, \code{unsafe},
\code{opaque}, \code{extern}, \code{native\_decide}, \code{sorry}, custom
\code{axiom}, and related escape hatches;
\item absence of the named scalar \code{contender\_10} from executable
\code{Contender11.lean} syntax;
\item separation of definition and proof imports;
\item presence of all stated definitions and theorems;
\item balanced ASCII delimiters in all new Lean files;
\item an exact strict target for the required comparison.
\end{itemize}

This is useful evidence against packaging mistakes, but it is not a Lean parser
or kernel.

\subsection{Independent finite reference model}

A small Python model reimplements only the structural algebra.  It does not
attempt System~$F_\omega$ normalization.  The test suite passed 670 assertions:

\begin{itemize}
\item 600 dense tests of
$\operatorname{unpair}_2(2^i(2m+1))=(i,m)$ for
$0\le i<12$ and $0\le m<50$;
\item child and input inequalities for a finite toy reflector;
\item strictness of every crown edge for inputs $0$ through $5$;
\item the exact seed--leaf--depth-zero-join--index comparison route;
\item the uniform toy predecessor bound;
\item a symbolic count of 43 crown edges at input 42.
\end{itemize}

The reference model is deliberately small enough to evaluate.  It tests the
shape of the proof and guards against off-by-one errors; it cannot certify the
Lean source or the inherited $\Fk$ metatheory.

\subsection{Outstanding kernel checks}

The environment used to produce this package did not contain \code{lake}.
Consequently the generated validation receipt says \code{not\_run}, rather than
fabricating a pass.  The included \code{check\_in\_repo.py} performs the
following when run in a pinned checkout:

\begin{enumerate}
\item build \code{BusyLean.Comparison10} as the accepted predecessor;
\item time \code{BusyLean.Contender11} against the 15-second definition budget;
\item time \code{BusyLean.Comparison11} against the 60-second proof budget;
\item run non-evaluating definitional smoke tests;
\item run \code{\#print axioms} for the new roots;
\item traverse transitive constant dependencies and reject earlier named
scalar contenders.
\end{enumerate}

A successful run of those commands is required before calling this an accepted
BusyLean contender.

%=====================================================================
\section{Integration and likely elaboration points}

The proposal is additive.  Copy the two modules to \code{BusyLean/}, append
their imports to \code{BusyLean.lean}, merge the audit fragment, and update the
checker so Contender~10 is the unbudgeted predecessor and Contender~11 is the
new definition phase.  The package's installer defaults to a dry run, prints
unified diffs, checks the pinned Git commit, and refuses to overwrite differing
candidate files.

The source closely mirrors already compiled repository patterns, but three
places deserve attention during the first kernel run.

\begin{enumerate}
\item \textbf{Root definitional equality.}  The proof of
\code{root\_program\_lt\_F11} uses \code{change} after unfolding one crown node
and one forest node.  If elaboration is more conservative than expected, two
explicit rewriting steps should replace \code{change}; the theorem statement
and construction need not change.
\item \textbf{Natural successor normal forms.}  The crown induction uses both
\code{k + 1} and successor induction.  These are definitionally aligned in the
expected Lean reduction, but a local \code{simpa [Nat.succ\_eq\_add\_one]} is a
harmless fallback.
\item \textbf{Timing attribution.}  A full historical clean build includes all
older normalization developments.  The contest rule budgets the new module on
top of a precompiled accepted predecessor.  The supplied checker follows that
three-phase discipline exactly.
\end{enumerate}

No proposed repair should introduce evaluation of a contender, oracle
monotonicity, new axioms, or a dependence on the named scalar.

%=====================================================================
\section{Potential objections and responses}

\subsection{``Is the crown merely padding?''}

The crown supplies a transparent diagonal margin, but the entry does not rest
on it.  Even before the crown, \code{family10\_lt\_stage11\_index} shows that
kind-polymorphic reflection along the self-generated index beats the whole
predecessor family.  The crown then makes the number of additional semantic
levels explicit and input-dependent.

\subsection{``Could Contender 10 write the same computation as one program?''}

A bounded $\Fk$ term at one reflection node may perform rich arithmetic and
many oracle calls.  The new construction nevertheless encloses every listed
such program in a new maximum, changes the oracle to a family of recursively
computed child stages, and repeats this operation along a meta-level $W$-tree.
The exact root-program and parent-edge theorems certify that structural
separation without requiring a semantic non-definability theorem for all
possible encodings.

\subsection{``Does the old W-language index limit the advance?''}

Yes, in a precise sense: the cofinality theorem ranges over trees produced by
the old W-language, because the current kind-polymorphic object language is a
numeral language.  This is a limitation, not a flaw in the comparison.  A
future object language with both kind polymorphism and positive inductive
$W$-types could generate a stronger self-index, but its normalization and
large-elimination metatheory would be a much larger formalization project.

\subsection{``Where is the exact proof-theoretic ordinal?''}

It is not asserted.  Kind polymorphism, the precise checker, the oracle
extension, forest reflection, and diagonal crown form a concrete computable
hierarchy.  Calibrating that hierarchy by a standard ordinal notation system
would be valuable separate research.  The BusyLean contest only needs the
formal finite inequality, which is stronger and less ambiguous than an
unproved ordinal analogy.

%=====================================================================
\section{Further research program}

\subsection{Immediate certification}

The first task is mechanical but essential: run Lean~4.34.0, repair any
elaboration-only issues, record definition and proof times, and inspect the
exact axiom list and dependency closure.  The package separates generic lemmas
from the final scalar so that repairs should not invite numerical reduction.

\subsection{A kind-polymorphic W-language}

The most direct qualitative successor would internalize positive inductive
families or $W$-types into the kind-polymorphic language, with a normalizer and
checker strong enough to support large elimination without inconsistency.
Then the self-generated index could enumerate trees produced by the same
language that performs reflection, rather than by an older predicative
language.  The central research problem is to design a total calculus that is
stronger than $\Fk$ yet still admits a reducibility or logical-relations proof
within the permitted Lean axiom envelope.

\subsection{Typed family oracles}

The current pairing exposes a family through a unary natural oracle.  A richer
object language could include a primitive family oracle or dependent oracle
interface, eliminating coding overhead and allowing typed child data.  To be a
contest advance rather than syntactic sugar, one would need an exact embedding
of the old interface and a program or tree not representable at bounded cost in
the old syntax.

\subsection{Self-generated kind indices}

Contender~10 can quantify over closed kinds but does not hand a syntactic kind
index back to the forest layer.  One could study a two-sorted hierarchy whose
nodes carry both $W$-tree branches and closed kind codes, with reflection over
all programs and type operators below a joint budget.  A diagonal index could
then traverse kind order and tree order simultaneously.

\subsection{Certified majorants and monotone reflection}

The present proof wisely avoids oracle monotonicity.  Another route is to
associate each accepted program with a certified monotone majorant.  Reflection
could then compare entire oracle families extensionally and support more
powerful transfer theorems.  The challenge is obtaining useful majorants for
conditionals and encoded data without collapsing the object language to an
obviously weaker positive fragment.

\subsection{Exact calibration}

A long-term mathematical project is to identify theories whose provably total
functions match the bounded maxima of the current $\Fk$ calculus, its oracle
extension, and the forest closure.  Even upper and lower ordinal bounds would
clarify which next language feature---large elimination, universes, inductive
families, or stronger recursion---produces a genuine proof-theoretic jump.

\subsection{Enumeration compression}

The current syntactic enumerations grow extremely quickly and use the same
input as both syntax depth and type-checking fuel.  Certified binary encodings,
hash-consed syntax, or direct rank enumerators might reduce checking overhead
without changing the mathematical maximum.  Any optimization must prove
coverage and semantic agreement and must preserve the contest's finite-time
build budget.

%=====================================================================
\section{Conclusion}

The proposed Contender~11 is
\[
  \boxed{\ContEleven=\Feleven(\Feleven(42))},
  \qquad
  \Feleven(n)=
  \Sstage_{\FamTen}(\Crown_{n+1}(I),n),
  \qquad
  \FamTen(n)=F_{10}(F_{10}(n)).
\]
It combines three mechanisms already justified separately in BusyLean:
kind-polymorphic normalization, family packing, and well-founded forest
reflection.  The new contribution is their composition at every forest node,
together with an input-diagonal crown and an exact proof architecture that
needs neither oracle monotonicity nor ordinal calibration.

The mathematical comparison is concise but strong:
\[
  \ContTen=\FamTen(42)
  <\Feleven(42)
  <\Feleven(\Feleven(42))=\ContEleven.
\]
The first inequality is backed by a leaf--join--self-index path and a strict
crown; the second is the general input bound.  Additional theorems expose the
program and tree margins.  The package is intentionally honest about its last
open engineering step: the Lean files must still be compiled and audited under
the pinned toolchain.  All material needed for that run is included.

%=====================================================================
\appendix
\section{Complete proposed definition module}
\label{app:def}

\lstinputlisting[caption={\code{BusyLean/Contender11.lean}}]{../candidate/BusyLean/Contender11.lean}

\section{Complete proposed comparison module}
\label{app:proof}

\lstinputlisting[caption={\code{BusyLean/Comparison11.lean}}]{../candidate/BusyLean/Comparison11.lean}

\section{Reproduction checklist}

\begin{enumerate}
\item Check out BusyLean at \pin.
\item Run \code{python3 integration/install.py /path/to/BusyLean} and review the
printed diff.
\item Re-run with \code{--apply}.
\item Run \code{python3 scripts/check\_in\_repo.py /path/to/BusyLean}.
\item Require successful definition/proof phases within 15 and 60 seconds.
\item Inspect \code{evidence/lean\_validation.external.json} and the full audit
output.
\item Rebuild this article with three \code{pdflatex} passes from the
\code{article/} directory.
\item Run \code{python3 scripts/static\_audit.py} and
\code{python3 scripts/reference\_model.py}.
\item Verify \code{SHA256SUMS}.
\end{enumerate}

\section{Artifact map}

\begin{longtable}{>{\raggedright\arraybackslash}p{0.45\linewidth}>{\raggedright\arraybackslash}p{0.47\linewidth}}
\toprule
Path & Purpose \\
\midrule
\endfirsthead
\toprule
Path & Purpose \\
\midrule
\endhead
\path{candidate/BusyLean/Contender11.lean} & Computational definitions. \\
\path{candidate/BusyLean/Comparison11.lean} & Comparison and structural proofs. \\
\path{candidate/test/Audit11.lean} & Axiom and transitive dependency audit. \\
\path{candidate/test/Contender11Smoke.lean} & Non-evaluating definitional checks. \\
\path{integration/install.py} & Pinned, dry-run-first installer. \\
\path{scripts/check_in_repo.py} & Kernel, timing, and audit runner. \\
\path{scripts/static_audit.py} & Source hygiene audit. \\
\path{scripts/reference_model.py} & Independent finite structural model. \\
\path{evidence/} & Machine-readable validation receipts and build logs. \\
\path{docs/VALIDATION_BOUNDARY.md} & Explicit trust boundary. \\
\path{article/contender11.tex} & This source. \\
\path{article/contender11.pdf} & Compiled report. \\
\bottomrule
\end{longtable}

\begin{thebibliography}{10}
\small

\bibitem{busylean}
BusyLean contributors.
\emph{BusyLean: name the biggest number, in Lean}.
Pinned revision \pin, October 2026.
\url{https://github.com/VladimirReshetnikov/BusyLean/tree/4c4b03d0e2003c479e3bee90a21c8bea13f66c40}.

\bibitem{roux}
Cody Roux and contributors.
\emph{name-the-biggest-number}.
\url{https://github.com/codyroux/name-the-biggest-number}.

\bibitem{aaronson}
Scott Aaronson.
\emph{Who Can Name the Bigger Number?}
\url{https://www.scottaaronson.com/writings/bignumbers.html}.

\bibitem{girard}
Jean-Yves Girard.
\emph{Interpr\'etation fonctionnelle et \'elimination des coupures de
l'arithm\'etique d'ordre sup\'erieur}.
Th\`ese d'\'Etat, Universit\'e Paris VII, 1972.

\bibitem{reynolds}
John C. Reynolds.
Towards a theory of type structure.
In \emph{Programming Symposium}, LNCS 19, pages 408--425, 1974.

\bibitem{barendregt}
Henk Barendregt.
Lambda calculi with types.
In Samson Abramsky, Dov Gabbay, and Thomas Maibaum, editors,
\emph{Handbook of Logic in Computer Science}, volume 2, pages 117--309.
Oxford University Press, 1992.

\bibitem{hurkens}
Antonius J. C. Hurkens.
A simplification of Girard's paradox.
In \emph{Typed Lambda Calculi and Applications}, LNCS 902, pages 266--278,
1995.

\bibitem{leanrec}
The Lean development team.
\emph{The Lean Language Reference: Recursive Definitions}.
\url{https://lean-lang.org/doc/reference/latest/Definitions/Recursive-Definitions/}.
Consulted October 2026; the package itself pins Lean~4.34.0.

\end{thebibliography}

\end{document}
